\chapter{Conclusion}\label{sec:Conclusion}

This work developed and validated a model-based dual estimation framework for the simultaneous 
estimation of the \ac{SOC} and \ac{SOH} of the WeLion SHP270-30 \ac{NMC} cell used in the eRacing 
electric motorcycle at ETSEIB. The framework combines experimental cell characterisation, 
algebraic parameter identification, and dual \ac{EKF} implementation in MATLAB. 
The following sections present the conclusions from each stage of the work.

\section{Battery Characterisation}

The discharge test confirmed an actual cell capacity of $27.616 \text{Ah}$, approximately $8\%$ below 
the nominal rating of $30 \text{Ah}$. This difference is consistent with typical behaviour in 
practice, where the rated nominal capacity represents an ideal cell. Factors during the discharge test 
such as the selected maximum and minimum voltage, conservative test conditions and ageing contributed to this difference.

The charge test confirmed the expected \ac{CC-CV} behaviour of the \ac{Li-ion} battery and established 
an energy efficiency of $91.5\%$, consistent with values reported for \ac{Li-ion} cells in 
the literature~\cite{ma16176063}. As expected, the Ah measured during charging exceeded the 
discharge capacity due to the inefficiencies; for this reason, the discharge test capacity was 
adopted as the reference for all \ac{SOC} calculations and simulations throughout this work.

The pulse discharge test produced an \ac{OCV}--\ac{SOC} curve spanning from $4.138\,\text{V}$ corresponding to a  
$100\%$ \ac{SOC} to $3.029\,\text{V}$ for $0\%$ \ac{SOC}. The characteristic nonlinear shape, with a
steeper slope at the higher and lower values of the \ac{SOC} and an almost linear region in between, 
is consistent with the typical behaviour of a \ac{NMC} lithium-ion cells~\cite{HassanTantaNMC}. 
The resulting lookup table was successfully implemented as the \ac{OCV}($z$) function in the 
equivalent-circuit model and in both the \ac{SOC} and \ac{SOH} \ac{EKF} algorithms.

\section{Parameter Identification}

The algebraic parameter identification method of Devarakonda and Hu~\cite{DEVARAKONDA2014} was 
applied to a two-level circuit. Twelve identification cycles were 
completed, covering a \ac{SOC} range from approximately $8\%$ to $99\%$.

The ohmic resistance $R_0$ and diffusion resistance $R_2$ were robustly identified across the full 
\ac{SOC} range, showing low variability between cycles. Their values, in the range of 
$2.1$--$2.6\,\text{m}\Omega$ and below $1.1\,\text{m}\Omega$ respectively, are consistent with 
expected values for a cell of this capacity~\cite{SOCSOH}.

The resistance $R_1$ and capacitance $C_1$ exhibited larger variability throughout the cycles, 
therefore two cycles were excluded by following the criteria ($|E_{\text{extracted}} - 
\text{OCV}(\text{SOC})| > 20\,\text{mV}$), and the validated cycles covered the range of 
$16.6\%$--$99.9\%$ \ac{SOC}. This behaviour is consistent with the sensitivity analysis described 
in Devarakonda and Hu~\cite{DEVARAKONDA2014}, when the test phase duration is short relative to 
the slow RC time constant, the algorithm cannot reliably separate the contributions of $R_1$ and 
$C_1$ from the initial capacitor voltage, making the individual values unreliable even when the 
overall voltage fit is reasonable.

The second-order Thévenin \acs{ECM} captures the dominant voltage response of the 
cell adequately for the purpose of state estimation, as confirmed by the \ac{EKF} 
validation results. However, the model cannot fully represent the complex 
electrochemical processes occurring in a cell. This trade-off between accuracy and 
computational feasibility is a recognised characteristic of equivalent-circuit 
models when compared to electrochemical or analytical approaches, 
which offer higher-fidelity representations of internal cell phenomena 
at the cost of implementability in real-time or dynamic simulation 
environments~\cite{DiazGonzalez2018}. 
In particular, the steep regions of the 
\ac{OCV}--\ac{SOC} curve, the fully charged and fully discharged phases. The 
complex behaviour in these regions cannot fully be captured by the linear RC 
structure. This is reflected in the numerical unreliability of $R_1$ and $C_1$ at 
low \ac{SOC} (approximately below $16\%$), and represents a known limitation of 
the \ac{ECM} approach at the extremes of the operating range.
\section{SOC Estimation with EKF}

The \ac{SOC} estimator was evaluated across three current profiles representative of different scenarios. 
Table~\ref{tab:SOC_RMSE} summarises the \ac{RMSE} results with respect 
to the Coulomb-counting reference.

\begin{table}[H]
    \centering
    \setlength{\arrayrulewidth}{0.7pt}
    \renewcommand{\arraystretch}{1.3} 
    \begin{tabular}{|l|c|}
        \hline
        \textbf{Test profile} & \textbf{RMSE [\%]} \\
        \hline
        Random steps ($0\,\text{A}$, $10\,\text{A}$, $30\,\text{A}$) & $0.691$ \\
        \hline
        Charge cycle (up to $20\,\text{A}$)                    & $0.942$ \\
        \hline
        Cyclic pulses ($10\,\text{A}$ and rest)                          & $1.141$ \\
        \hline
    \end{tabular}
    \caption{SOC estimation \ac{RMSE} for the \ac{EKF} \ac{SOC} filter across the three test profiles.}
    \label{tab:SOC_RMSE}
\end{table}

All three profiles produced \ac{RMSE} values below $1.15\%$, and the estimation error remained 
within the $3\sigma$ bounds throughout each test. These results are consistent with values reported 
in comparable \ac{EKF}--\ac{ECM} implementations in the literature, where errors on the order of 
$1\%$--$2\%$ are typical for similar profiles~\cite{EKFEV}\cite{BatteryKalmanRaceCar}.

The random step-current profile produced the lowest \ac{RMSE}. The charge test produced 
an intermediate error.

\section{Simultaneous SOC and SOH Estimation with Dual EKF}

The dual \ac{EKF} was applied to the same three profiles, simultaneously estimating \ac{SOC} and 
the \ac{SOH} parameters $R_0$ and $Q$. Table~\ref{tab:DualEKF_results} summarises the \ac{SOC} 
\ac{RMSE} and the final \ac{SOH} estimates for each test.

\begin{table}[H]
    \centering
    \small
    \setlength{\arrayrulewidth}{0.7pt} 
    \renewcommand{\arraystretch}{1.3} 
    \begin{tabular}{|l|c|c|c|}
        \hline
        \textbf{Profile} & \textbf{SOC RMSE [\%]} & $Q$ [Ah] & $R_0$ [m$\Omega$] \\
        \hline
        Random steps  & $0.958$ & $27.92$  & $2.132$ \\
        \hline
        Charge cycle  & $1.390$ & $28.44$  & $2.246$ \\
        \hline
        Cyclic pulses & $0.943$ & $27.67$  & $2.134$ \\
        \hline
    \end{tabular}
    \caption{Dual EKF results: SOC RMSE and final SOH estimates across the three test profiles.}
    \label{tab:DualEKF_results}
\end{table}

% The \ac{SOC} \ac{RMSE} of the dual filter was higher than that of the \ac{SOC}
% filter in all three profiles, which is expected, since during the parameter convergence transient the dual 
% filter uses a slightly incorrect $Q$ to propagate the \ac{SOC}, introducing a transient error not 
% present when $Q$ is fixed. Once the parameter filter converges, the \ac{SOC} error stabilises and 
% remains within bounds.

The capacity estimate $Q$ converged to values close to the characterised reference of 
$27.616\,\text{Ah}$ for the discharge profiles, with the tightest convergence observed for the 
cyclic pulse test $27.67\,\text{Ah}$, $3\sigma$. The charge cycle 
produced a slightly higher estimate ($28.44\,\text{Ah}$), which reflects the Coulombic 
inefficiency of the charging process. The total energy transferred during a full charge 
exceeds the energy provided in a discharge cycle, and therefore the parameter filter partially interprets 
additional charge as a larger capacity.

The resistance estimates $R_0$ in all three profiles converged to values consistent with 
those in the characterisation table at the corresponding \ac{SOC} levels. Values of 
$R_0$ exceeding $100\%$ indicate that the online estimate is lower than the nominal 
table value, reflecting that the filter adapts to the actual operating resistance, which can 
differ from the characterisation value due to differences in current magnitude, temperature, and 
operating history between the two tests.

Overall, these results demonstrate that the dual \ac{EKF} framework is capable of tracking 
both \ac{SOC} and the \ac{SOH} parameters simultaneously and consistently across different 
current profiles.

\section{Limitations and Future Work}

The equivalent-circuit model does not account for temperature effects nor hysterisis. All experiments were 
conducted at approximately ambient temperature but in a racing scenario, the cell temperature can vary 
significantly during a race, which affects all model parameters. Incorporating temperature-dependent 
parameter tables would improve the accuracy and robustness of the estimator 
under real operating conditions.

The parameter identification at low and high \ac{SOC} is incomplete due to the 
numerical sensitivity of the algebraic method in regions where the \ac{OCV} gradient is steep,
a methodology accounting for the electrochemical behaviour in these regions would improve the results.

Future work could also explore the \ac{UKF} as an alternative to the \ac{EKF}. The \ac{UKF} avoids 
the first-order linearisation used by the \ac{EKF}, and may provide improved accuracy in the 
nonlinear regions of the \ac{OCV}--\ac{SOC} curve without a significant increase in 
computational cost~\cite{plett2015battery}.

Finally, validation under actual race current profiles, including the high peak currents and 
rapid transients characteristic of a motorcycle competition, would further confirm the applicability of 
the framework to the operational conditions for which it was designed.